\section{Rational slopes}\label{sec:rational}
\subsection{Integer slopes for \texorpdfstring{$m\ge4$}{m >= 4}}
\begin{proposition}\label{prop:integer}
Let $m$ and $k$ be integers and let $M=m-2$. Assume the following.
\begin{hypotheses}
\item\label{hyp:integer-arrows} $m\ge4$.
\item\label{hyp:integer-range} $k\in[M]_{\Z}$.
\end{hypotheses}
Let
\[
 \mathcal S_m=
 \left\{(a,b)\in\Z^2:\
 \begin{gathered}
 b>a\ge2,\quad \gcd(a,b)=1,\\
 |R|_N\in[M]_{\Z}\text{ for every simple run }R\subseteq\mathcal U(a,b)
 \end{gathered}\right\}.
\]
Then, we have the following.

\begin{enumerate}
\item\label{c:integer-limit} We have
\[
 \lim_{\substack{a\to\infty,\ (a,b)\in\cA_m\cap\mathcal S_m\\\frac{b}{a}\to k}}
 \frac{\log E_m(a,b)}a=\ell(R_k).
\]
\item\label{c:integer-endpoints} The following one-sided limits hold.
\begin{enumerate}[label=(\alph*),ref=\theenumi(\alph*)]
\item\label{c:integer-left-endpoint} We have
\[
 \lim_{\substack{a\to\infty,\ (a,b)\in\cA_m\\\frac ba\downarrow1}}
 \frac{\log E_m(a,b)}a=\ell(R_1).
\]
\item\label{c:integer-right-endpoint} We have
\[
 \lim_{\substack{a\to\infty,\ (a,b)\in\cA_m\\\frac ba\uparrow M}}
 \frac{\log E_m(a,b)}a=\ell(R_M).
\]
\end{enumerate}
\end{enumerate}
\end{proposition}
\begin{proof}
\prf{Claim~\ref{c:integer-limit}} Let us estimate the number of vertical run-edits needed to obtain a constant run.
Consider $(a,b)\in\cA_m\cap\mathcal S_m$ as $a\to\infty$ and $\frac ba\to k$.
For all sufficiently large $a$, the first $a$ simple runs of $\mathcal{U}(a,b)$ have
vertical lengths $k-1$ and $k$ or $k$ and $k+1$. The $(a+1)$-st simple run of $\mathcal U(a,b)$ is $R_1$. The total vertical run-edit distance to $R_k^{a+1}$ is
\[
 \Delta=|b-1-ka|+(k-1)=o(a).
\]

Let us prove the upper bound. Let $\nu=a+1$ and $P=(m+2)\Delta+1$. We change the run $\mathcal{U}(a,b)$ to the constant run $R_k^\nu$ by vertical run-edits. Every added outer simple run has vertical length $k$. Hence, the target is $R_k^{\nu+2P+2}$.
Assumptions~\ref{hyp:integer-arrows}--\ref{hyp:integer-range} and $(a,b)\in\mathcal S_m$ give the length and insertion hypotheses for Claim~\ref{c:composition-bound} of Proposition~\ref{prop:composition}.
Equations~\eqref{eq:composition} and~\eqref{eq:dictionary} give
\[
 mE_m(a,b)\le \operatorname{Fib}_m(\nu,D)\,C_0(R_k^{\nu+2P+2}),\qquad D\le\Delta.
\]

Let us check the error term as $a\to\infty$ and $\frac ba\to k$ through $\cA_m\cap\mathcal S_m$.
Set $x=\frac{\Delta}{\nu}$. We have $x\to0$ and $\frac\nu a\to1$.
Since $x\log x\to0$ as $x\downarrow0$, Definition~\ref{def:entropy-loss-function} gives $\psi_m(x)\to0$.
Equation~\eqref{eq:entropy-loss}, applied with $N=\nu$, gives
\[
 0\le\frac{\log\operatorname{Fib}_m(\nu,D)}a
 \le\frac\nu a\psi_m\left(\frac\Delta\nu\right)\longrightarrow0.
\]
Also, $P=o(\nu)$ gives $\frac{\nu+2P+2}{a}\to1$ and $\nu+2P+2\to\infty$.
Claim~\ref{c:periodic-limit} of Proposition~\ref{prop:periodic} therefore gives
\[
\begin{aligned}
 \frac{\log E_m(a,b)}a
 &\le\frac{\log \operatorname{Fib}_m(\nu,D)}a
 +\frac{\nu+2P+2}{a}\frac{\log C_0(R_k^{\nu+2P+2})}{\nu+2P+2}
 -\frac{\log m}{a},\\
 \limsup_{\substack{a\to\infty,\ (a,b)\in\cA_m\cap\mathcal S_m\\\frac{b}{a}\to k}}
 \frac{\log E_m(a,b)}a&\le\ell(R_k).
\end{aligned}
\]

Let us prove the lower bound. Let $A_1,\ldots,A_\nu$ be simple runs such that $\mathcal U(a,b)=A_1\ccn \cdots\ccn A_\nu$.
Let $l=\nu-2P-2$. Since $\frac P\nu\to0$, we have $l\ge1$ for all sufficiently large $a$.
We start from $R_k^l$. For the current path $Z$,
we use the run-edits $Z\mapsto\operatorname{Ins}(Z,i,N)$ at $i\in[|Z|]_{\Z}$
or $Z\mapsto\operatorname{Del}(Z,v)$ at a vertical edge $v$.
At most $\Delta$ such run-edits give a translated copy of $A_{P+2}\ccn \cdots\ccn A_{\nu-P-1}$.
We use $A_1,\ldots,A_{P+1}$ and $A_{\nu-P},\ldots,A_\nu$ for the outer simple runs. Equations~\eqref{eq:composition} and~\eqref{eq:dictionary} give
\[
 mE_m(a,b)\ge\frac{C_0(R_k^l)}{\operatorname{Fib}_m(l,D)},\qquad D\le\Delta.
\]
Since $\frac{l}{a}\to1$ and $D=o(l)$, Claim~\ref{c:periodic-limit} of Proposition~\ref{prop:periodic} and~\eqref{eq:entropy-loss} give
\[
\begin{aligned}
 \frac{\log E_m(a,b)}a
 &\ge\frac{l}{a}\frac{\log C_0(R_k^l)}l
 -\frac{\log \operatorname{Fib}_m(l,D)}a-\frac{\log m}a,\\
 \liminf_{\substack{a\to\infty,\ (a,b)\in\cA_m\cap\mathcal S_m\\\frac{b}{a}\to k}}
 \frac{\log E_m(a,b)}a&\ge\ell(R_k).
\end{aligned}
\]

\prf{Claim~\ref{c:integer-left-endpoint}} Let us prove the one-sided limit as $\frac ba\downarrow1$. We have $1<\frac ba<M$ for all sufficiently large $a$.
Since $b-a\ge1$, equation~\eqref{eq:ambient} gives $1<\Gamma(a,b)<M$.
The floor differences in~\eqref{eq:frame-word} lie in $[M]_{\Z}$.
The $(a+1)$-st simple run of $\mathcal U(a,b)$ is $R_1$.
Hence, $(a,b)\in\mathcal S_m$.
Thus, Claim~\ref{c:integer-limit} applies with $k=1$.

\prf{Claim~\ref{c:integer-right-endpoint}} Let us prove the one-sided limit as $\frac ba\uparrow M$. We have $1<\Gamma(a,b)<\frac ba<M$ for all sufficiently large $a$ by~\eqref{eq:ambient}.
The floor differences in~\eqref{eq:frame-word} and the appended $R_1$ have vertical lengths in $[M]_{\Z}$.
Hence, $(a,b)\in\mathcal S_m$. Claim~\ref{c:integer-limit} applies with $k=M$.
\end{proof}

\Needspace{16\baselineskip}
\subsection{Mechanical run sequences near a noninteger rational slope}
\begin{definition}\label{def:mechanical-paths}
Let $m\in\Z_{\ge3}$, $p\in\Z_{>0}$, and $q\in\Z_{>0}$ satisfy $p\ge q$ and $\gcd(p,q)=1$.
\begin{enumerate}[font=\itshape]
\item\label{d:rational-period} Define
\[
 M_{\mathrm{run}}(m)=\begin{cases}2,&m=3,\\m-2,&m\ge4.\end{cases}
\]
For $i\in\Z_{>0}$, define
\[
 w_i(p,q)=\floor{\frac{ip}{q}}-\floor{\frac{(i-1)p}{q}}.
\]
We define the \textit{rational run} and the \textit{rational run sequence} by
\[
 W_{p,q}=R_{w_1(p,q)}\ccn\cdots\ccn R_{w_q(p,q)},\qquad
 \mathcal R_{p,q}=(R_{w_i(p,q)})_{i\in\Z_{>0}}.
\]
For $j\in\Z_{>0}$ and $\alpha=\frac pq\in(1,M_{\mathrm{run}}(m))$, define
\begin{equation}\label{eq:rational-period}
 A_j(p,q)=C_0(W_{p,q}^j),\qquad \Upsilon_m(\alpha)=\frac{\ell(W_{p,q})}q.
\end{equation}
\item\label{d:mechanical-sequence} Let $s\ge1$ and $u\in[0,1)$ be real numbers. Define
\[
 \operatorname{MechRun}(s,u)=(R_{j_i})_{i\in\Z_{>0}},\qquad
 j_i=\floor{u+is}-\floor{u+(i-1)s}.
\]
We call $\operatorname{MechRun}(s,u)$ the
\textit{mechanical run sequence} of slope $s$ and intercept $u$.
\item Let $\alpha=\frac pq$.
For $h\in\Z_{\ge0}$, define
\[
 \operatorname{RunPhase}_q(h)=h-q\floor{\frac hq}.
\]
We call $\operatorname{RunPhase}_q(h)$ the \textit{run phase} at the start of the $(h+1)$-st simple run of $\mathcal R_{p,q}$.
\item Let $n\in\Z_{>0}$ and let $A_1,\ldots,A_n$ be simple runs.
Define $W=A_1\ccn \cdots\ccn A_n$.
Let $r\in\Z_{>0}$ and $\ell\in\Z_{>0}$ satisfy $r+\ell-1\le n$.
Define the \textit{$(r,\ell)$-run-factor} of $W$ by
\[
 \operatorname{RunFactor}_{r,\ell}(W)
 =W\left[\sum_{i=1}^{r-1}|A_i|,\ \sum_{i=1}^{r+\ell-1}|A_i|\right].
\]

Let $\mathcal A=(A_i)_{i\in\Z_{>0}}$ be a sequence of simple runs.
For $r\in\Z_{>0}$ and $\ell\in\Z_{>0}$, define
\[
 \operatorname{RunFactor}_{r,\ell}(\mathcal A)
 =\operatorname{RunFactor}_{r,\ell}(A_1\ccn \cdots\ccn A_{r+\ell-1}).
\]
\end{enumerate}
\end{definition}

\begin{lemma}\label{lem:rational-mechanical}
Let $m\in\Z_{\ge3}$.
Let $p\in\Z_{>0}$ and $q\in\Z_{>0}$ satisfy $\gcd(p,q)=1$, $q\ge2$, and
$1<\frac pq<M_{\mathrm{run}}(m)$.
Let $\alpha=\frac pq$.
Let $s\ge1$ and $u\in[0,1)$. Assume the following.
\begin{hypotheses}
\item\label{hyp:rational-mechanical-slope} $|s-\alpha|<q^{-2}$.
\end{hypotheses}
Then, we have the following.

\begin{enumerate}
\item\label{c:rational-mechanical-short} For every $r\in\Z_{>0}$ and $l\in[q-1]_{\Z}$, there exists $t\in\Z_{>0}$ such that the path
\[
 \operatorname{RunFactor}_{r,l}(\operatorname{MechRun}(s,u))
\]
is a translated copy of $\operatorname{RunFactor}_{t,l}(\mathcal R_{p,q})$.
\item\label{c:rational-mechanical-period} For $r\in\Z_{>0}$, let
$Z_r=\operatorname{RunFactor}_{r,q}(\operatorname{MechRun}(s,u))$.
If $|Z_r|_N=p$, then $Z_r$ is a translated copy of a cyclic rotation
of $W_{p,q}$ by simple runs.
\item\label{c:rational-mechanical-lengths} Let $k=\floor\alpha$.
Then, $w_i(p,q)\in[k,k+1]_{\Z}$ for $i\in[q]_{\Z}$ and $w_1(p,q)=k$.
\item\label{c:rational-mechanical-zero-phase} For $i\in[0,q-1]_{\Z}$, we have
\[
 \floor{is}=\floor{i\alpha}.
\]
\end{enumerate}
\end{lemma}
\begin{proof}
\prf{Claim~\ref{c:rational-mechanical-short}} Let us locate the floor thresholds within one period.
For $i\in[q-1]_{\Z}$, define
$T_i(t)=\ceil{\frac{ip}{q}}-it$ for $t\in\R$.
Let us prove
\begin{equation}\label{eq:mechanical-threshold-order}
 T_i(s)<T_j(s)\quad\Longleftrightarrow\quad T_i(\alpha)<T_j(\alpha)
 \qquad(i,j\in[q-1]_{\Z}).
\end{equation}
Since $\alpha=\frac pq$, the numbers $T_i(\alpha)$
are distinct elements of $\{\frac1q,\ldots,\frac{q-1}{q}\}$.
By assumption~\ref{hyp:rational-mechanical-slope},
\[
 |T_i(s)-T_i(\alpha)|=i|s-\alpha|<\frac1q,
 \qquad 0<T_i(s)<1.
\]
If $i\in[q-1]_{\Z}$ and $j\in[q-1]_{\Z}$ are distinct, then we also have
\[
 \left|(T_i(s)-T_j(s))-(T_i(\alpha)-T_j(\alpha))\right|
 =|i-j|\,|s-\alpha|<\frac1q\le|T_i(\alpha)-T_j(\alpha)|.
\]
Hence, equation~\eqref{eq:mechanical-threshold-order} holds.
For $t\in\{s,\alpha\}$ and $v\in[0,1)$,
\[
 \floor{v+it}=
 \begin{cases}
 \ceil{i\alpha}-1,&v<T_i(t),\\
 \ceil{i\alpha},&v\ge T_i(t).
 \end{cases}
\]
Thus, equation~\eqref{eq:mechanical-threshold-order} gives, for every $l\in[q-1]_{\Z}$,
\begin{equation}\label{eq:mechanical-prefix-vectors}
 \left\{(\floor{v+is})_{i\in[l]_{\Z}}:v\in[0,1)\right\}
 =\left\{(\floor{v+i\alpha})_{i\in[l]_{\Z}}:v\in[0,1)\right\}.
\end{equation}

Let us apply~\eqref{eq:mechanical-prefix-vectors} to $\operatorname{RunFactor}_{r,l}(\operatorname{MechRun}(s,u))$.
Let $v_r=u+(r-1)s-\floor{u+(r-1)s}$.
For $i\in[l]_{\Z}$,
\begin{equation}\label{eq:mechanical-factor-prefix}
\begin{aligned}
 \left|\operatorname{RunFactor}_{r,i}(\operatorname{MechRun}(s,u))\right|_N
 &=\floor{u+(r+i-1)s}-\floor{u+(r-1)s}\\
 &=\floor{v_r+is}.
\end{aligned}
\end{equation}
Every vector in the right-hand side of~\eqref{eq:mechanical-prefix-vectors} occurs at an intercept in
$\{0,\frac1q,\ldots,\frac{q-1}{q}\}$.
Since $\gcd(p,q)=1$,
\[
 \left\{(t-1)\alpha-\floor{(t-1)\alpha}:t\in[q]_{\Z}\right\}
 =\left\{0,\frac1q,\ldots,\frac{q-1}{q}\right\}.
\]
Equations~\eqref{eq:mechanical-prefix-vectors} and~\eqref{eq:mechanical-factor-prefix} therefore give
$t\in[q]_{\Z}$ such that
\[
 \left|\operatorname{RunFactor}_{r,i}(\operatorname{MechRun}(s,u))\right|_N
 =\left|\operatorname{RunFactor}_{t,i}(\mathcal R_{p,q})\right|_N
 \qquad(i\in[l]_{\Z}).
\]
Taking consecutive differences shows that
$\operatorname{RunFactor}_{r,l}(\operatorname{MechRun}(s,u))$
is a translated copy of $\operatorname{RunFactor}_{t,l}(\mathcal R_{p,q})$.

\prf{Claim~\ref{c:rational-mechanical-period}} Let us determine $Z_r$ from its first $q-1$ simple runs and its total vertical length.
Claim~\ref{c:rational-mechanical-short} gives $t\in\Z_{>0}$ such that the first $q-1$ simple runs
of $Z_r$ are a translated copy of $\operatorname{RunFactor}_{t,q-1}(\mathcal R_{p,q})$.
Since $q\alpha=p\in\Z$, we have
\[
 \left|\operatorname{RunFactor}_{t,q}(\mathcal R_{p,q})\right|_N
 =\floor{(t+q-1)\alpha}-\floor{(t-1)\alpha}=p.
\]
The assumption $|Z_r|_N=p$ therefore determines the last simple-run length.
Thus, $Z_r$ is a translated copy of $\operatorname{RunFactor}_{t,q}(\mathcal R_{p,q})$.
The identity $\floor{z+p}=\floor z+p$ for $z\in\R$ shows that the simple-run lengths of $\mathcal R_{p,q}$
are $q$-periodic. Hence, $\operatorname{RunFactor}_{t,q}(\mathcal R_{p,q})$ is a cyclic rotation of $W_{p,q}$ up to translation.

\prf{Claim~\ref{c:rational-mechanical-lengths}} Let us compute the periodic run lengths.
Let $\beta=\alpha-k\in[0,1)$.
For $i\in[q]_{\Z}$, we have
\[
 w_i(p,q)=k+\floor{i\beta}-\floor{(i-1)\beta}\in[k,k+1]_{\Z}.
\]
At $i=1$, we have $w_1(p,q)=k+\floor\beta-\floor0=k$.

\prf{Claim~\ref{c:rational-mechanical-zero-phase}} Let us bound the distance of $i\alpha$ from the integers.
For $i=0$, both floors are zero.
For $i\in[q-1]_{\Z}$, coprimality of $p$ and $q$ gives
\[
 \floor{i\alpha}+\frac1q\le i\alpha
 \le\floor{i\alpha}+1-\frac1q.
\]
Assumption~\ref{hyp:rational-mechanical-slope} gives
\[
 |is-i\alpha|<\frac{i}{q^2}<\frac1q.
\]
Hence, $\floor{i\alpha}<is<\floor{i\alpha}+1$, so $\floor{is}=\floor{i\alpha}$.
\end{proof}

\begin{definition}\label{sync:rational:264}
Let $m\in\Z_{\ge3}$, $p\in\Z_{>0}$, and $q\in\Z_{>0}$ satisfy $\gcd(p,q)=1$ and $q\ge2$.
Let $\alpha=\frac pq\in(1,M_{\mathrm{run}}(m))$.
Let $n\in\Z_{>0}$ and let $A_1,\ldots,A_n$ be simple runs.
Define $W=A_1\ccn \cdots\ccn A_n$.
For $r\in[n-q+1]_{\Z}$, define
\[
 Z_r=\operatorname{RunFactor}_{r,q}(W).
\]
We call $Z_r$ \textit{$\alpha$-bad} if $|Z_r|_N\ne p$.
Define
\[
 D_{\mathrm{bad}}(W,p,q)=\#\{r\in[n-q+1]_{\Z}:|Z_r|_N\ne p\}.
\]
\end{definition}

\begin{lemma}[Counting exceptional floor differences]\label{lem:floor-discrepancies}
Let $q\in\Z_{>0}$, $N\in\Z_{>0}$, $p\in\Z$, $s\in\R$, and $u\in\R$.
Assume $|qs-p|<1$.
Then,
\[
 \#\{i\in[0,N-1]_{\Z}:\floor{u+(i+q)s}-\floor{u+is}\ne p\}
 \le N|qs-p|+q.
\]
\end{lemma}
\begin{proof}
Let us express the exceptional floors as jumps of an integer sequence.
Let $\epsilon=qs-p$ and define
\[
 \delta_i=\floor{u+(i+q)s}-\floor{u+is}-p
 \qquad(i\in[0,N-1]_{\Z}).
\]
Since $p\in\Z$ and $|\epsilon|<1$,
\[
 \delta_i=\floor{u+is+\epsilon}-\floor{u+is}\in
 \begin{cases}
 [0,1]_{\Z},&\epsilon\ge0,\\
 [-1,0]_{\Z},&\epsilon<0.
 \end{cases}
\]
Thus, $\delta_i\ne0$ implies $\delta_i=1$ when $\epsilon\ge0$ and $\delta_i=-1$ when $\epsilon<0$.

Let us sum the floor differences in each residue class.
For $r\in[0,q-1]_{\Z}$, let
$n_r=\#\{i\in[0,N-1]_{\Z}:i\equiv r\pmod q\}$.
Let us compute $\sum_{t=0}^{n_r-1}\delta_{r+tq}$.
Telescoping gives
\[
\begin{aligned}
 \sum_{t=0}^{n_r-1}\delta_{r+tq}
 &=\floor{u+(r+n_rq)s}-\floor{u+rs}-n_rp,\\
 \left|\sum_{t=0}^{n_r-1}\delta_{r+tq}-n_r\epsilon\right|&<1.
\end{aligned}
\]
Since $\delta_i\in\{0,1\}$ for $\epsilon\ge0$ and $\delta_i\in\{-1,0\}$ for $\epsilon<0$, we obtain
\[
\begin{aligned}
 \#\{i\in[0,N-1]_{\Z}:\delta_i\ne0\}
 &=\sum_{r=0}^{q-1}\left|\sum_{t=0}^{n_r-1}\delta_{r+tq}\right|\\
 &\le\sum_{r=0}^{q-1}(n_r|\epsilon|+1)=N|\epsilon|+q.
\end{aligned}
\]
\end{proof}

\begin{lemma}\label{lem:bad-factors}
Let $m\in\Z_{\ge3}$, $p\in\Z_{>0}$, and $q\in\Z_{>0}$ satisfy $\gcd(p,q)=1$ and $q\ge2$.
Let $\alpha=\frac pq\in(1,M_{\mathrm{run}}(m))$.
Let $b>a\ge2$ be coprime integers.
Then, we have the following.
\begin{enumerate}
\item\label{c:bad-factors-bound} Assume the following.
\begin{hypotheses}
\item\label{hyp:bad-factor-length} $a\ge q$.
\item\label{hyp:bad-factor-slope} $|\Gamma(a,b)-\alpha|<q^{-2}$.
\end{hypotheses}
Then,
\begin{equation}\label{eq:bad-windows}
 D_{\mathrm{bad}}(\mathcal U(a,b),p,q)\le |qb-pa|+q+2.
\end{equation}
\item\label{c:bad-factors-negligible} We have
\[
 D_{\mathrm{bad}}(\mathcal U(a,b),p,q)=o(a)
 \qquad\left(a\to\infty,\ \frac ba\to\alpha,\ b>a,\ \gcd(a,b)=1\right).
\]
\end{enumerate}
\end{lemma}
\begin{proof}
\prf{Claim~\ref{c:bad-factors-bound}} Let us bound the number of $\alpha$-bad run-factors contained in the first $a$ simple runs of $\mathcal U(a,b)$.
Assumptions~\ref{hyp:bad-factor-length} and~\ref{hyp:bad-factor-slope} give $a-q+1\ge1$ and $|q\Gamma(a,b)-p|<1$.
Lemma~\ref{lem:floor-discrepancies} with $N=a-q+1$, $s=\Gamma(a,b)$, and $u=0$ gives
\begin{equation}\label{eq:frame-floor-exceptions}
\begin{aligned}
 &\#\left\{i\in[0,a-q]_{\Z}:\floor{(i+q)\Gamma(a,b)}-\floor{i\Gamma(a,b)}\ne p\right\}\\
 &\qquad\le(a-q+1)|q\Gamma(a,b)-p|+q.
\end{aligned}
\end{equation}
For $i\in[0,a-q]_{\Z}$, Definition~\ref{def:frame} gives
\[
 \left|\operatorname{RunFactor}_{i+1,q}(\mathcal U(a,b))\right|_N
 =\floor{(i+q)\Gamma(a,b)}-\floor{i\Gamma(a,b)}.
\]
By~\eqref{eq:frame-word}, the $(a+1)$-st simple run of $\mathcal U(a,b)$ is $R_1$.
For $r\in[a-q+2]_{\Z}$, the path $\operatorname{RunFactor}_{r,q}(\mathcal U(a,b))$
contains the $(a+1)$-st simple run if and only if $r=a-q+2$.
Hence, equation~\eqref{eq:frame-floor-exceptions} gives
\[
\begin{aligned}
 D_{\mathrm{bad}}(\mathcal U(a,b),p,q)
 &\le 1+(a-q+1)|q\Gamma(a,b)-p|+q\\
 &\le aq|\Gamma(a,b)-\alpha|+q+1.
\end{aligned}
\]
Assumption~\ref{hyp:bad-factor-length} and~\eqref{eq:ambient} give $\mu_1(a,b)\ge a+1>q$ and
\[
 aq|\Gamma(a,b)-\alpha|
 =\left|qb-pa-\frac q{\mu_1(a,b)}\right|
 \le |qb-pa|+\frac q{\mu_1(a,b)}.
\]
Since $\frac q{\mu_1(a,b)}<1$, inequality~\eqref{eq:bad-windows} follows.

\prf{Claim~\ref{c:bad-factors-negligible}} Let us verify the range of the ambient slope.
Equation~\eqref{eq:ambient} gives
\[
 |\Gamma(a,b)-\alpha|\le\left|\frac ba-\alpha\right|+\frac1{a(a+1)}\longrightarrow0.
\]
For all sufficiently large $a$, Claim~\ref{c:bad-factors-bound} applies and gives
\[
 0\le\frac{D_{\mathrm{bad}}(\mathcal U(a,b),p,q)}a
 \le\left|q\frac ba-p\right|+\frac{q+2}a\longrightarrow0.
\]
\end{proof}

\begin{lemma}\label{lem:separated-subwords}
Let $m\in\Z_{\ge3}$, $p\in\Z_{>0}$, and $q\in\Z_{>0}$ satisfy $\gcd(p,q)=1$ and $q\ge2$.
Let $\alpha=\frac pq\in(1,M_{\mathrm{run}}(m))$.

\begin{enumerate}
\item\label{c:separated-subwords-rotations} For every integer $J\ge4$, each cyclic rotation $V$ of $W_{p,q}^J$
by simple runs satisfies $C_0(V)\ge A_{J-3}(p,q)$.
\item\label{c:separated-subwords-construction} Let $h\ge1$ be an integer.
Let $X_1,\ldots,X_h$ be runs.
Let $B_0,\ldots,B_h$ each consist of three simple runs. Assume the following.
\begin{hypotheses}
\item\label{hyp:separated-lengths} Every simple run of $B_0,\ldots,B_h$ has vertical length in $[M_{\mathrm{run}}(m)]_{\Z}$.
\end{hypotheses}
Then,
\begin{equation}\label{eq:separated-subwords}
 C_0(B_0\ccn X_1\ccn B_1\ccn \cdots\ccn X_h\ccn B_h)\ge\prod_{i=1}^h C_0(X_i).
\end{equation}
\end{enumerate}
\end{lemma}
\begin{proof}
\prf{Claim~\ref{c:separated-subwords-rotations}} Let us apply Lemma~\ref{lem:outer-shell-extension} to a rotated periodic path.
Let $t=J-3$ and let $j\in[0,q-1]_{\Z}$.
Let $U_j=\operatorname{CRot}(W_{p,q},j)$.
Let $P=2q-j$ and $S=q+j$.
Let $X=\operatorname{RunFactor}_{P+1,qt}(U_j^{t+3})$. Then,
\[
 X=\operatorname{Trans}(W_{p,q}^t,\operatorname{Start}_{\mathrm{pt}}(X)),\qquad
 P+qt+S=q(t+3).
\]
Since $P+j=2q$, we have $\operatorname{RunPhase}_q(P+j)=0$.
Also, $\max(P,S)\le2\min(P,S)$.
Claim~\ref{c:rational-mechanical-lengths} of Lemma~\ref{lem:rational-mechanical} gives
$|R|_N\in[M_{\mathrm{run}}(m)]_{\Z}$ for every simple run $R$ of $W_{p,q}$.
Lemma~\ref{lem:outer-shell-extension} gives
\[
 C_0(U_j^{t+3})\ge C_0(W_{p,q}^t)=A_{J-3}(p,q).
\]

\prf{Claim~\ref{c:separated-subwords-construction}} Let us specify the shell paths.
Let $B_i=R_{b_{i,1}}\ccn R_{b_{i,2}}\ccn R_{b_{i,3}}$ for $i\in[0,h]_{\Z}$.
For $r\in[h]_{\Z}$ and $s\in[h]_{\Z}$ satisfying $r\le s$ and
$(u,v)\in\{(2,2),(2,1),(1,3),(3,2)\}$, define
\[
 T_{uv}(r,s)=
 \left(\mathop{\ccn}\limits_{j=4-u}^{3}R_{b_{r-1,j}}\right)
 \ccn X_r\ccn\left(\mathop{\ccn}\limits_{i=r+1}^{s}\bigl(B_{i-1}\ccn X_i\bigr)\right)
 \ccn\left(\mathop{\ccn}\limits_{j=1}^{v}R_{b_{s,j}}\right).
\]
The iterated concatenations use increasing index order.
We compare paths up to translation using Lemma~\ref{lem:word-reindexing}.
For $r=s$, we have
\[
\begin{aligned}
 T_{22}(r,r)&=\operatorname{Shell}
 (\operatorname{Shell}(X_r;b_{r-1,3},b_{r,1},0);b_{r-1,2},b_{r,2},0),\\
 T_{21}(r,r)&=\operatorname{Shell}(R_{b_{r-1,3}},X_r;b_{r-1,2},b_{r,1},0),\\
 T_{13}(r,r)&=\operatorname{Shell}(X_r,R_{b_{r,1}}\ccn R_{b_{r,2}};b_{r-1,3},b_{r,3},0),\\
 T_{32}(r,r)&=\operatorname{Shell}(T_{21}(r,r);b_{r-1,1},b_{r,2},0).
\end{aligned}
\]
For $r<s$, we have
\[
\begin{aligned}
 T_{22}(r,s)&=\operatorname{Shell}
 (\operatorname{Shell}(X_r;b_{r-1,3},b_{r,1},0),T_{21}(r+1,s);b_{r-1,2},b_{s,2},0),\\
 T_{21}(r,s)&=\operatorname{Shell}(T_{13}(r,s-1),X_s;b_{r-1,2},b_{s,1},0),\\
 T_{13}(r,s)&=\operatorname{Shell}(X_r,T_{32}(r+1,s);b_{r-1,3},b_{s,3},0),\\
 T_{32}(r,s)&=\operatorname{Shell}
 (T_{22}(r,s-1),\operatorname{Shell}(X_s;b_{s-1,3},b_{s,1},0);b_{r-1,1},b_{s,2},0).
\end{aligned}
\]

Let us count the resulting colored gradings.
Assumption~\ref{hyp:separated-lengths} gives $b_{i,j}\in[M_{\mathrm{run}}(m)]_{\Z}\subseteq[m-1]_{\Z}$.
Lemma~\ref{lem:word-bounds} gives $C_0(R_{b_{i,j}})\ge1$ and
$C_0(R_{b_{i,1}}\ccn R_{b_{i,2}})\ge1$.
Claim~\ref{c:shell-recolor} of Lemma~\ref{lem:shell} gives an injective map
into colored gradings of surplus zero at each shell step.
Induction on $s-r$ therefore gives
\[
 C_0(T_{uv}(r,s))\ge\prod_{i=r}^{s}C_0(X_i).
\]
Finally,
\[
 \operatorname{Shell}(T_{22}(1,h);b_{0,1},b_{h,3},0)
 =B_0\ccn X_1\ccn B_1\ccn \cdots\ccn X_h\ccn B_h
\]
up to translation. Claim~\ref{c:shell-recolor} of Lemma~\ref{lem:shell}
and Lemma~\ref{lem:word-reindexing} give~\eqref{eq:separated-subwords}.
\end{proof}

\begin{theorem}[Noninteger rational all-sequence limit]\label{thm:rational}
Let $m\in\Z_{\ge3}$, $p\in\Z_{>0}$, and $q\in\Z_{>0}$ satisfy $\gcd(p,q)=1$ and $q\ge2$.
Let $\alpha=\frac pq\in(1,M_{\mathrm{run}}(m))$.
Then,
\begin{equation}\label{eq:rational-limit}
 \lim_{\substack{a\to\infty,\ (a,b)\in\cA_m\\\frac{b}{a}\to\alpha}}
 \frac{\log E_m(a,b)}a=\frac{\ell(W_{p,q})}q.
\end{equation}
\end{theorem}
\begin{proof}
Let us construct the initial factor and the separated periodic factors.
Let $k=\floor\alpha$.
In the limit in~\eqref{eq:rational-limit}, we have
$a\ge q$ and $|\Gamma(a,b)-\alpha|<q^{-2}$ for all sufficiently large $a$.
\prf{the upper bound}
First, we divide $\mathcal{U}(a,b)$ into periodic pieces.
Equation~\eqref{eq:frame-word} and Claims~\ref{c:rational-mechanical-lengths} and~\ref{c:rational-mechanical-zero-phase}
of Lemma~\ref{lem:rational-mechanical} give
\[
 \operatorname{RunFactor}_{1,1}(\mathcal U(a,b))
 =\operatorname{RunFactor}_{1,1}(W_{p,q})=R_k,\qquad
 \operatorname{RunPhase}_q(0)=0.
\]
Let $0=c_0<c_1<\cdots<c_d=a$ satisfy
\[
 \{c_1,\ldots,c_{d-1}\}
 =\{r\in[a-q+1]_{\Z}:|\operatorname{RunFactor}_{r,q}(\mathcal U(a,b))|_N\ne p\}.
\]
For $i\in[d]_{\Z}$, let
\[
 X_i=\operatorname{RunFactor}_{c_{i-1}+1,c_i-c_{i-1}}(\mathcal U(a,b)).
\]
Each $X_i$ contains no $\alpha$-bad run-factor. We have
$d\le D_{\mathrm{bad}}(\mathcal U(a,b),p,q)+1$.
By Claim~\ref{c:rational-mechanical-short} of Lemma~\ref{lem:rational-mechanical}, a piece with $l<q$ simple runs
is a translated copy of $\operatorname{RunFactor}_{t,l}(\mathcal R_{p,q})$ for some $t\in\Z_{>0}$.
In a piece with at least $q$ simple runs, every run-factor consisting of $q$ simple runs has
vertical length $p$. Subtracting the sums for consecutive run-factors gives
\[
 0=\sum_{i=j+1}^{j+q}\Lambda_i(a,b)-\sum_{i=j}^{j+q-1}\Lambda_i(a,b)
 =\Lambda_{j+q}(a,b)-\Lambda_j(a,b).
\]
Thus, $\Lambda_{j+q}(a,b)=\Lambda_j(a,b)$ for $i\in[d]_{\Z}$ and $c_{i-1}<j<j+q\le c_i$.
Claim~\ref{c:rational-mechanical-period} of Lemma~\ref{lem:rational-mechanical} identifies the first $q$ simple runs of each such piece with a cyclic rotation of $W_{p,q}$.

Second, let us align each piece with a run-factor of $\mathcal R_{p,q}$.
Let us prove that no insertion is needed before the first piece.
Claim~\ref{c:rational-mechanical-zero-phase} of Lemma~\ref{lem:rational-mechanical}, applied with $s=\Gamma(a,b)$, gives
\[
 \floor{i\Gamma(a,b)}=\floor{i\alpha}\qquad(i\in[q-1]_{\Z}).
\]
If the first piece has at least $q$ simple runs, its first $q$ lengths sum to $p$.
Thus, its $q$-th length also agrees with $w_q(p,q)$.
The identities $\Lambda_{j+q}(a,b)=\Lambda_j(a,b)$ for $1\le j<j+q\le c_1$ identify $X_1$ with $\operatorname{RunFactor}_{1,c_1}(\mathcal R_{p,q})$.
Consequently, no insertion is needed before $X_1$.

Let us proceed by induction through $X_2,\ldots,X_d$.
For the current piece $X$, choose the least $t\in[q]_{\Z}$ such that
$X$ is a translated copy of $\operatorname{RunFactor}_{t,|X|_E}(\mathcal R_{p,q})$.
Let $P\ccn X\ccn Z$ be the current path and let $h=|P|_E$.
The induction hypothesis is that the simple-run lengths of $P$ are $w_1(p,q),\ldots,w_h(p,q)$.
The identification of $X_1$ gives the base case. Define
\[
 j=\operatorname{RunPhase}_q\bigl(q+t-1-\operatorname{RunPhase}_q(h)\bigr)
 \in[0,q-1]_{\Z}.
\]
Then, $\operatorname{RunPhase}_q(h+j)=t-1$.
If $j>0$, insert the simple runs of $\operatorname{RunFactor}_{h+1,j}(\mathcal R_{p,q})$
after $P$, producing
\[
 \operatorname{Ins}\bigl(P\ccn X\ccn Z,|P|,\operatorname{RunFactor}_{h+1,j}(\mathcal R_{p,q})\bigr).
\]
If $j=0$, the path is $P\ccn X\ccn Z$.

By $q$-periodicity, the simple-run lengths of $X$ now agree with those at positions $h+j+1,\ldots,h+j+|X|_E$ in $\mathcal R_{p,q}$.
Thus, the prefix through $X$ has simple-run lengths $w_1(p,q),\ldots,w_{h+j+|X|_E}(p,q)$, which proves the induction step.

Let us complete the final period before $R_1$.
Let $h$ be the number of simple runs before the final $R_1$ of $\mathcal U(a,b)$ after these insertions.
Let $j=\operatorname{RunPhase}_q\bigl(q-\operatorname{RunPhase}_q(h)\bigr)$.
If $j>0$, insert $\operatorname{RunFactor}_{h+1,j}(\mathcal R_{p,q})$ before $R_1$.
Then, $\operatorname{RunPhase}_q(h+j)=0$.
Let $K$ be the number of inserted simple runs. Let $n_0=\frac{a+K}{q}$.
The result is
\[
 W_{p,q}^{n_0}\ccn R_1,\qquad
 K\le(q-1)(D_{\mathrm{bad}}(\mathcal U(a,b),p,q)+2).
\]
For each insertion, use~\eqref{eq:run-insertion-map},
\eqref{eq:run-insertion-grades}, and~\eqref{eq:run-insertion-recolor}.
The resulting grading $g$ satisfies $\operatorname{Surp}(W_{p,q}^{n_0}\ccn R_1,g)=K$.

Third, let us add the outer simple runs.
Apply Lemma~\ref{lem:outer-shell-extension} with $(P,S)=(q,q-1)$
to add $W_{p,q}$ before $W_{p,q}^{n_0}\ccn R_1$ and $R_{w_2(p,q)}\ccn \cdots\ccn R_{w_q(p,q)}$ after $W_{p,q}^{n_0}\ccn R_1$.
The target is
\begin{equation}\label{eq:rational-initial-extension}
 Q_0=W_{p,q}^{n_0+1}\ccn R_1\ccn R_{w_2(p,q)}\ccn \cdots\ccn R_{w_q(p,q)}.
\end{equation}
If $k=1$, we have $Q_0=W_{p,q}^{n_0+2}$. For $m=3$, the condition $1<\alpha<2$ gives $k=1$.

Let us obtain the initial run of length $k$.
For $k>1$, let us replace the factor $R_1$ in~\eqref{eq:rational-initial-extension} by $R_k$.
We have $m\ge4$. Apply Claims~\ref{c:vertical-insert-connected}--\ref{c:vertical-insert-injective} of Lemma~\ref{lem:vertical-insert} $k-1$ times.
For each application, use $q$ unary shell paths to add one copy of $W_{p,q}$ on each side.
After the $j$th vertical-edge insertion, the path is
\[
 Q_j=W_{p,q}^{n_0+j+1}\ccn R_{j+1}\ccn R_{w_2(p,q)}\ccn \cdots\ccn R_{w_q(p,q)}\ccn W_{p,q}^j
 \qquad(j\in[k-1]_{\Z}).
\]
The transported shell grading $h_j$ satisfies $\operatorname{Surp}(Q_j,h_j)=K$.

Let us reduce the surplus to zero.
We have $Q_{k-1}=W_{p,q}^{n_0+2k}$.
Let $H=\ceil{\frac Kq}$.
Apply Lemma~\ref{lem:surplus-reduction} with $Q=W_{p,q}^{n_0+2k}$, $n=qH$, and
\[
 a_j=w_{q-\operatorname{RunPhase}_q(j-1)}(p,q),\qquad
 b_j=w_{1+\operatorname{RunPhase}_q(j-1)}(p,q)
 \qquad(j\in[qH]_{\Z}).
\]
The final path is a translated copy of $W_{p,q}^{n_0+2k+2H}$.

Let us prove injectivity of the composition.
At each simple-run insertion, equations~\eqref{eq:run-insertion-grades} and~\eqref{eq:run-insertion-recolor} recover the input grading and coloring.
At each vertical-edge insertion, Claim~\ref{c:vertical-insert-injective} of Lemma~\ref{lem:vertical-insert} gives injectivity.
Lemma~\ref{lem:outer-shell-extension} gives injectivity of each outer extension.
Lemma~\ref{lem:word-reindexing} and the injection~\eqref{eq:surplus-reduction-map} give
\begin{equation}\label{eq:rational-upper}
 mE_m(a,b)\le A_N(p,q),\qquad
 N=\frac{a+K}q+2k+2\ceil{\frac{K}{q}}.
\end{equation}

Let us obtain the upper limit from~\eqref{eq:rational-upper}.
By Claim~\ref{c:bad-factors-negligible} of Lemma~\ref{lem:bad-factors}, we have $K=o(a)$ with $q$ fixed. Hence, $\frac{N}{a}\to\frac1q$. Claim~\ref{c:periodic-limit} of Proposition~\ref{prop:periodic} gives
\[
 \begin{aligned}
 \frac{\log E_m(a,b)}a&\le\frac{N}{a}\frac{\log A_N(p,q)}{N}-\frac{\log m}{a},\\
 \limsup_{\substack{a\to\infty,\ (a,b)\in\cA_m\\\frac{b}{a}\to\alpha}}
 \frac{\log E_m(a,b)}a&\le\frac{\ell(W_{p,q})}q.
 \end{aligned}
\]

\prf{the lower bound}
Let us select separated subpaths on which to count gradings.
Let $J\ge4$ be an integer and let
\[
 h=\floor{\frac{a-2}{Jq+3}}
\]
if $a$ is sufficiently large and $h\ge1$. Define
\[
\begin{aligned}
 X_i&=\operatorname{RunFactor}_{4+(i-1)(Jq+3),Jq}(\mathcal U(a,b))
       &&(i\in[h-1]_{\Z}),\\
 X_h&=\operatorname{RunFactor}_{4+(h-1)(Jq+3),\,a-5-(h-1)(Jq+3)}(\mathcal U(a,b)),\\
 B_i&=\operatorname{RunFactor}_{i(Jq+3)+1,3}(\mathcal U(a,b))
       &&(i\in[0,h-1]_{\Z}),\\
 B_h&=\operatorname{RunFactor}_{a-1,3}(\mathcal U(a,b)).
\end{aligned}
\]
Then, $\mathcal U(a,b)=B_0\ccn X_1\ccn B_1\ccn \cdots\ccn X_h\ccn B_h$ and $|X_h|_E\ge Jq$.

Let us count colored gradings on these subpaths.
Let
\[
 \mathcal I=\{h\}\cup\{i\in[h-1]_{\Z}:D_{\mathrm{bad}}(X_i,p,q)>0\}.
\]
Then, $\#\mathcal I\le D_{\mathrm{bad}}(\mathcal U(a,b),p,q)+1$.
We call $X_i$ exceptional for $i\in\mathcal I$.
For $i\in[h]_{\Z}\setminus\mathcal I$, the path $X_i$ is a translated copy of $\operatorname{RunFactor}_{t,Jq}(\mathcal R_{p,q})$ for some $t\in\Z_{>0}$. By Claim~\ref{c:separated-subwords-rotations} of Lemma~\ref{lem:separated-subwords}, such a subpath has at least $A_{J-3}(p,q)$ connected colored gradings of surplus zero.
Each exceptional subpath $X$ has a connected colored grading $(g,\gamma)$ with
$\operatorname{Surp}(X,g)=0$ by Lemma~\ref{lem:word-bounds}.
Claim~\ref{c:separated-subwords-construction} of Lemma~\ref{lem:separated-subwords} gives
\[
 mE_m(a,b)\ge A_{J-3}(p,q)^{\max(0,h-D_{\mathrm{bad}}(\mathcal U(a,b),p,q)-1)}.
\]

Let us take the lower limit and then let $J$ tend to infinity.
If $J$ and $q$ are fixed, then~\eqref{eq:bad-windows} gives
\[
 \frac{h}{a}\longrightarrow\frac1{Jq+3},\qquad
 \frac{D_{\mathrm{bad}}(\mathcal U(a,b),p,q)+1}{a}\longrightarrow0
 \quad\left(a\to\infty,\ \frac ba\to\alpha,\ (a,b)\in\cA_m\right).
\]
We then let $J\to\infty$ with $q$ fixed. By Claim~\ref{c:periodic-limit} of Proposition~\ref{prop:periodic}, the lower bound becomes
\[
 \liminf_{\substack{a\to\infty,\ (a,b)\in\cA_m\\\frac{b}{a}\to\alpha}}
 \frac{\log E_m(a,b)}a
 \ge\lim_{J\to\infty}\frac{\log A_{J-3}(p,q)}{Jq+3}
 =\frac{\ell(W_{p,q})}q.
\]
\end{proof}

